\section{Reflections}
Reflecting on the project at hand, a couple of learnings stand out. When we started refactoring the legacy version of the application, we transitioned to a Golang code-base. Picking a language unfamiliar to the majority of us made the refactoring process pretty steep. Refactoring the unit-tests of the legacy app, for example, demanded the use of Go-libraries that functioned vastly different than frameworks such as pytest (if one tries to run the code from the initial commit of test-files, 8bbbaff, it is clearly broken). This learning-by-doing approach, however, made the evolution of the program easier, as the time spent understanding the new code base, made any evolutionary steps more fluent.

The logging and monitoring parts of the operation process was one of the main differences to other software projects we have worked on. The simulation really forced us to actually prioritize backlog items for deployment in a more intentional matter, as a large amount of user's 'depended' on our service. In a software side-project, you can develop new features as you find them interesting, but with 100.000 users, an unoptimized http-handler quickly becomes a top priority. This also made maintenance more of a matter of urgency than previous projects we have worked on. Something breaking has consequences for a vast amount of users, and sometimes, that demanded late night debugging. All of this was only possible by our monitoring and logging frameworks, and the course has given us great insights into working with such in the future. 